\section{Preference Learning: Detailed Proofs}
\label{app:preference-proofs}

This appendix provides complete proofs for the preference learning equivalence and gap theorems.

\subsection{Proof Strategy Overview}

The preference learning proofs follow a unified architecture across DPO, GRPO-PL, and GRPO-RL:

\paragraph{The Lipschitz Chain.}
Every gap bound is established by composing Lipschitz constants through a three-stage chain:
\[
\text{Loss} \xrightarrow{L_1\text{-Lip}} \text{Logit/Score} \xrightarrow{L_2\text{-Lip}} \text{Policy ratios} \xrightarrow{L_3\text{-Lip}} \text{Oracle value}
\]

The final Lipschitz constant is the product: $L = L_1 \cdot L_2 \cdot L_3$.

\paragraph{Oracle-Measurability as Factorization.}
Oracle-measurable policies satisfy $\pi(\cdot \mid x) = \pi(\cdot \mid \fstar(x))$. This allows rewriting any oracle-measurable quantity as a function of $\fstar(x)$ rather than $x$ directly. When $\fstar(Z) = \fstar(X)$, all such quantities are identical.

\paragraph{From Pointwise to Expected.}
Pointwise Lipschitz bounds integrate to expected bounds:
\[
|f(x) - f(x')| \le L \cdot d(x, x') \implies |\E_X[f(X)] - \E_{X'}[f(X')]| \le L \cdot \E[d(X, X')].
\]
This is the key step connecting sample-level analysis to population gap bounds.

\subsection{Oracle-Measurability Foundations}

\begin{definition}[Oracle-Measurable Policy]
A policy $\pi: \Strings \to \Delta(\A)$ is \emph{oracle-measurable} if:
\[
\metric(\fstar(x), \fstar(x')) = 0 \implies \pi(\cdot \mid x) = \pi(\cdot \mid x')
\]
\end{definition}

\begin{definition}[Oracle-Indexed Preference]
A preference function $\text{pref}: \Strings \times \A \times \A \to [0, 1]$ is \emph{oracle-indexed} if:
\[
\metric(\fstar(x), \fstar(x')) = 0 \implies \text{pref}(x, a, a') = \text{pref}(x', a, a')
\]
for all actions $a, a'$.
\end{definition}

\subsection{DPO Loss Structure}

\begin{definition}[Bradley-Terry Preference Model]
\[
P(a \succ a' \mid x) = \sigma(\beta \cdot (s(x, a) - s(x, a')))
\]
where $s(x, a) = \log \frac{\pi_\theta(a \mid x)}{\pi_{\text{ref}}(a \mid x)}$ and $\sigma(z) = \frac{1}{1 + e^{-z}}$.
\end{definition}

\begin{definition}[DPO Loss]
\[
\mathcal{L}_{\text{DPO}}(\theta; x, a_w, a_l) = -\log \sigma\left(\beta \cdot \left(\log \frac{\pi_\theta(a_w \mid x)}{\pi_{\text{ref}}(a_w \mid x)} - \log \frac{\pi_\theta(a_l \mid x)}{\pi_{\text{ref}}(a_l \mid x)}\right)\right)
\]
\end{definition}

\subsection{Lipschitz Properties}

\begin{lemma}[Sigmoid Lipschitz]\label{lem:sigmoid-lip}
The sigmoid function $\sigma(z) = \frac{1}{1+e^{-z}}$ is $\frac{1}{4}$-Lipschitz:
\[
|\sigma(z) - \sigma(z')| \le \frac{1}{4}|z - z'|
\]
\end{lemma}

\begin{proof}
$\sigma'(z) = \sigma(z)(1 - \sigma(z))$. Since $\sigma(z) \in (0, 1)$, we have $\sigma'(z) \le \frac{1}{4}$ with maximum at $z = 0$.
\end{proof}

\begin{lemma}[Negative Log-Sigmoid Lipschitz]\label{lem:neglog-lip}
The function $\ell(z) = -\log \sigma(z)$ is $1$-Lipschitz for $z \in \mathbb{R}$:
\[
|\ell(z) - \ell(z')| \le |z - z'|
\]
\end{lemma}

\begin{proof}
$\ell'(z) = \sigma(-z) \in (0, 1)$, so $|\ell'(z)| < 1$.
\end{proof}

\begin{assumption}[Policy-Lipschitz]\label{assum:policy-lip}
The log-ratio difference is Lipschitz in the oracle value:
\[
\left|\log \frac{\pi_\theta(a \mid x)}{\pi_{\text{ref}}(a \mid x)} - \log \frac{\pi_\theta(a \mid x')}{\pi_{\text{ref}}(a \mid x')}\right| \le L_{\text{pol}} \cdot \metric(\fstar(x), \fstar(x'))
\]
for oracle-measurable policies.
\end{assumption}

\begin{theorem}[DPO Pointwise Lipschitz]\label{thm:dpo-lip}
Under Assumption~\ref{assum:policy-lip}:
\[
|\mathcal{L}_{\text{DPO}}(\theta; x, a_w, a_l) - \mathcal{L}_{\text{DPO}}(\theta; x', a_w, a_l)| \le 2|\beta| L_{\text{pol}} \cdot \metric(\fstar(x), \fstar(x'))
\]
\end{theorem}

\begin{proof}
Let $s(x,a)=\log\frac{\pi_\theta(a\mid x)}{\pi_{\text{ref}}(a\mid x)}$ and
$\Lambda(x)=s(x,a_w)-s(x,a_l)$. The loss is
$\mathcal{L}_{\text{DPO}}(\theta;x,a_w,a_l)=-\log\sigma(\beta\,\Lambda(x))$.
By Lemma~\ref{lem:neglog-lip}, $-\log\sigma$ is $1$-Lipschitz, so
\[
|\mathcal{L}(x)-\mathcal{L}(x')|\le |\beta|\,|\Lambda(x)-\Lambda(x')|.
\]
By the triangle inequality,
$
|\Lambda(x)-\Lambda(x')|
\le |s(x,a_w)-s(x',a_w)|+|s(x,a_l)-s(x',a_l)|.
$
Assumption~\ref{assum:policy-lip} bounds each term by
$L_{\text{pol}}\metric(\fstar(x),\fstar(x'))$. Combining yields
\[
|\mathcal{L}(x)-\mathcal{L}(x')|
\le 2|\beta|L_{\text{pol}}\metric(\fstar(x),\fstar(x')),
\]
as claimed.
\end{proof}

\subsection{Training Equivalence}

\begin{theorem}[DPO Training Equivalence]\label{app:thm:dpo-equiv}
Assume:
\begin{enumerate}
    \item $\metric(\fstar(Z^{(R)}(X)), \fstar(X)) = 0$ almost surely
    \item The policy $\pi_\theta$, reference $\pi_{\text{ref}}$, and preference oracle are all oracle-measurable/indexed
\end{enumerate}
Then the sets of population minimizers coincide:
\[
\argmin_\theta \E_X[\mathcal{L}_{\text{DPO}}(\theta; X)] = \argmin_\theta \E_Z[\mathcal{L}_{\text{DPO}}(\theta; Z^{(R)}(X))]
\]
\end{theorem}

\begin{proof}
Let $Y=\fstar(X)$. Oracle-measurability of the policies and oracle-indexing of the pair sampler imply that, conditional on $Y$, the joint law of $(a_w,a_l)$ and the log-ratio terms is identical whether the input is $X$ or $Z^{(R)}(X)$. By the Doob--Dynkin lemma there exists $\tilde{\mathcal{L}}_{\text{DPO}}$ with
\[
\mathcal{L}_{\text{DPO}}(\theta; x, a_w, a_l) = \tilde{\mathcal{L}}_{\text{DPO}}(\theta; \fstar(x), a_w, a_l).
\]
Exact preservation gives $\fstar(Z^{(R)}(X))=Y$ a.s., hence
\[
\E_X[\mathcal{L}_{\text{DPO}}(\theta; X)]
= \E_Y\!\left[\E[\tilde{\mathcal{L}}_{\text{DPO}}(\theta; Y, a_w, a_l)\mid Y]\right]
= \E_Z[\mathcal{L}_{\text{DPO}}(\theta; Z^{(R)}(X))].
\]
Since $\mathcal{L}^X_{\text{DPO}}(\theta) = \mathcal{L}^Z_{\text{DPO}}(\theta)$ for all $\theta$, the argmin sets coincide.
\end{proof}

\subsection{Quantitative Gap Bounds}

\begin{theorem}[DPO Gap Bound]\label{app:thm:dpo-gap}
Let $\Delta_R = \E[\metric(\fstar(Z^{(R)}(X)), \fstar(X))]$. Then:
\[
|\mathcal{L}^X_{\text{DPO}}(\theta) - \mathcal{L}^Z_{\text{DPO}}(\theta)| \le 2|\beta| L_{\text{pol}} \cdot \Delta_R
\]
\end{theorem}

\begin{proof}
\begin{align*}
|\mathcal{L}^X_{\text{DPO}}(\theta) - \mathcal{L}^Z_{\text{DPO}}(\theta)|
&= |\E_{X,a,a'}[\mathcal{L}_{\text{DPO}}(\theta; X, a_w, a_l) - \mathcal{L}_{\text{DPO}}(\theta; Z^{(R)}(X), a_w, a_l)]| \\
&\le \E_{X,a,a'}[|\mathcal{L}_{\text{DPO}}(\theta; X, a_w, a_l) - \mathcal{L}_{\text{DPO}}(\theta; Z^{(R)}(X), a_w, a_l)|] \\
&\le 2|\beta| L_{\text{pol}} \cdot \E_X[\metric(\fstar(Z^{(R)}(X)), \fstar(X))] \\
&= 2|\beta| L_{\text{pol}} \cdot \Delta_R
\end{align*}
\end{proof}

\subsection{GRPO-PL Bounds}

\begin{definition}[Plackett-Luce Model]
For $k$ candidates with scores $s_1, \ldots, s_k$:
\[
P(\sigma \mid s_1, \ldots, s_k) = \prod_{i=1}^{k} \frac{\exp(s_{\sigma(i)})}{\sum_{j=i}^{k} \exp(s_{\sigma(j)})}
\]
\end{definition}

\begin{lemma}[Plackett-Luce Lipschitz]
The negative log Plackett-Luce probability is Lipschitz in the scores:
\[
\left|-\log P(\sigma \mid s) + \log P(\sigma \mid s')\right| \le (k-1) \|s - s'\|_\infty
\]
\end{lemma}

\begin{proof}
Expand the log-likelihood:
\[
-\log P(\sigma \mid s)=\sum_{i=1}^{k-1}\left(\log\sum_{j=i}^k e^{s_{\sigma(j)}}-s_{\sigma(i)}\right).
\]
The log-sum-exp map is $1$-Lipschitz in $\|\cdot\|_\infty$ because its gradient is a probability simplex vector. The linear term is also $1$-Lipschitz in $\|\cdot\|_\infty$. Therefore each summand is $1$-Lipschitz, and summing $k-1$ terms yields
\[
\big|-\log P(\sigma \mid s)+\log P(\sigma \mid s')\big|\le (k-1)\|s-s'\|_\infty.
\]
\end{proof}

\begin{theorem}[GRPO-PL Gap Bound]
Under oracle-measurable policies and group generators:
\[
|\mathcal{L}^X_{\text{GRPO-PL}}(\theta) - \mathcal{L}^Z_{\text{GRPO-PL}}(\theta)| \le L_{\text{grpo}} \cdot \Delta_R
\]
where $L_{\text{grpo}} = (k-1) \cdot L_{\text{pol}}$.
\end{theorem}

\begin{proof}
Let $s(X)$ denote the score vector with entries $s_i(X)=\log\pi(a_i\mid X)-\log\pi_{\text{ref}}(a_i\mid X)$. By Lemma~(Plackett-Luce Lipschitz),
\[
\big|-\log P(\sigma\mid s(X))+\log P(\sigma\mid s(Z))\big|
\le (k-1)\|s(X)-s(Z)\|_\infty.
\]
Oracle-measurability and the policy-Lipschitz assumption give
$\|s(X)-s(Z)\|_\infty \le L_{\text{pol}}\metric(\fstar(X),\fstar(Z))$.
Taking expectations over $(X,Z^{(R)}(X),\sigma)$ yields
$|\mathcal{L}^X_{\text{GRPO-PL}}(\theta)-\mathcal{L}^Z_{\text{GRPO-PL}}(\theta)| \le (k-1)L_{\text{pol}}\Delta_R$.
\end{proof}

\subsection{GRPO-RL Bounds}

\begin{definition}[GRPO-RL Objective]
\[
\mathcal{L}_{\text{GRPO-RL}}(\theta) = -\E\left[\frac{1}{k}\sum_{i=1}^k \min(\rho_i A_i, \text{clip}(\rho_i) A_i)\right] + \beta \cdot D_{\text{KL}}(\pi_\theta \| \pi_{\text{ref}})
\]
where $\rho_i = \pi_\theta(a_i \mid x) / \pi_{\text{old}}(a_i \mid x)$ and $A_i = (r_i - \bar{r})/\sigma_r$.
\end{definition}

\begin{assumption}[Reward-Lipschitz]
The reward function is Lipschitz in the oracle value:
\[
|r(x, a) - r(x', a)| \le L_{\text{rew}} \cdot \metric(\fstar(x), \fstar(x'))
\]
\end{assumption}

\begin{theorem}[GRPO-RL Gap Bound]
Under oracle-measurable policies and rewards, and assuming $\sigma_r \ge \sigma_{\min} > 0$ on sampled groups:
\[
|\mathcal{L}^X_{\text{GRPO-RL}}(\theta) - \mathcal{L}^Z_{\text{GRPO-RL}}(\theta)| \le L_{\text{grpo-rl}} \cdot \Delta_R
\]
where $L_{\text{grpo-rl}}$ depends on $L_{\text{pol}}$, $L_{\text{rew}}$, $\sigma_{\min}$, and the clipping parameter.
\end{theorem}

\begin{proof}
Assume rewards are bounded and the empirical standard deviation satisfies $\sigma_r\ge\sigma_{\min}>0$. The reward-Lipschitz condition implies the reward vector is Lipschitz in $\metric(\fstar(X),\fstar(Z))$, and the z-score map $r\mapsto(r-\bar r)/\sigma_r$ is Lipschitz on $\{\sigma_r\ge\sigma_{\min}\}$ with constant $L_A$ depending on $(L_{\text{rew}},\sigma_{\min})$. The clipping operator is 1-Lipschitz and the KL term is Lipschitz under the policy-Lipschitz assumption. Combining these bounds yields a constant $L_{\text{grpo-rl}}$ such that
$|\mathcal{L}^X_{\text{GRPO-RL}}(\theta)-\mathcal{L}^Z_{\text{GRPO-RL}}(\theta)|\le L_{\text{grpo-rl}}\Delta_R$.
\end{proof}

\subsection{Unified Framework}

\begin{theorem}[Unified Preference Gap]\label{app:thm:unified-gap}
For any expected preference learning loss with $L$-Lipschitz inner expectation:
\[
|\E_X[\mathcal{L}(\theta; X)] - \E_Z[\mathcal{L}(\theta; Z^{(R)}(X))]| \le L \cdot \Delta_R
\]
\end{theorem}

\begin{proof}
Let $(X,Z)$ denote the coupled pair $(X,Z^{(R)}(X))$. By assumption,
$|\mathcal{L}(\theta;X)-\mathcal{L}(\theta;Z)|\le L\,\metric(\fstar(X),\fstar(Z))$.
Taking expectations and applying Jensen's inequality gives
\[
|\E[\mathcal{L}(\theta;X)]-\E[\mathcal{L}(\theta;Z)]|
\le \E\big[|\mathcal{L}(\theta;X)-\mathcal{L}(\theta;Z)|\big]
\le L\,\E\big[\metric(\fstar(X),\fstar(Z))\big]
= L\,\Delta_R.
\]
\end{proof}

The adaptive-tree optimizer-transfer results are collected separately in
Appendix~\ref{app:adaptive-tree-proofs}, where we give full human-readable
proofs for both the expected-tree perturbation theorem and its
high-probability version.
